% !TeX root = ../偏微分方程笔记.tex
% 由主文档通过 \input 载入；请编译主文档。
\setcounter{chapter}{17} % 保留原书章号。
\chapter{数值方法与综合建模}
\label{chap:numerics}

数值方法的任务是以有限个未知量近似连续问题。一个计算结果可用，既需要离散方程反映原方程，也需要误差在迭代中受到控制。

\section{有限差分与局部截断误差}
在 $[0,\ell]$ 上取 $J$ 个等长小区间，$h=\ell/J$，时间步长为 $\tau$。记 $x_j=jh$、$t_m=m\tau$、$U_j^m\approx u(x_j,t_m)$，并定义中心差分 $D_hU_j=(U_{j+1}-2U_j+U_{j-1})/h^2$。一维热方程的显式\emph{有限差分法（finite difference method）}为
\begin{equation}
 U_j^{m+1}=\mu U_{j-1}^m+(1-2\mu)U_j^m+\mu U_{j+1}^m+\tau f(x_j,t_m),
 \qquad \mu=\frac{\kappa\tau}{h^2}.
 \label{eq:heat-explicit-scheme}
\end{equation}
输入是网格、系数、初值及边值；先置 $U_j^0=g(x_j)$，每个时间层赋边界值，再更新内部点，迭代到指定终止时刻。每步需要 $O(J)$ 运算；若只保存相邻时间层，存储为 $O(J)$。

将精确解代入差分方程并除以时间步长，所得残差称\emph{局部截断误差（local truncation error）}。若 $u_{tt}$ 与 $u_{xxxx}$ 等导数在所考虑闭区域有界，泰勒展开给出
\[
 \frac{u(x_j,t_{m+1})-u(x_j,t_m)}\tau-\kappa D_hu(x_j,t_m)-f(x_j,t_m)
 =O(\tau+h^2).
\]
这称\emph{一致性（consistency）}，尚未证明数值解接近精确解。

\section{稳定性、收敛与步长限制}
当 $0\le\mu\le1/2$ 时，式\eqref{eq:heat-explicit-scheme}的三个权重非负且和为一。对同样边界与源项的两组数值解之差，每步上确界范数不增，这就是相应的\emph{稳定性（stability）}。把精确解残差作为误差方程的源项迭代，若初始和边界值精确且 $m\tau\le T$，则
\[
 \max_j|U_j^m-u(x_j,t_m)|\le T\max_{j,r}|\text{残差}_{j,r}|
 \le C_T(\tau+h^2).
\]
因此在上述光滑性和步长限制下得到\emph{收敛性（convergence）}。这段论证把一致性、稳定性和收敛之间的作用分别写明。

对无限均匀或周期网格，代入空间频率 $e^{ij\theta}$，放大因子为 $G=1-4\mu\sin^2(\theta/2)$；要求 $|G|\le1$ 同样得到标准限制 $\mu\le1/2$。后向欧拉格式和克兰克--尼科尔森格式（Crank--Nicolson scheme）分别为
\[
 (I-\tau\kappa D_h)U^{m+1}=U^m+\tau f^{m+1},\qquad
 (I-\tfrac\tau2\kappa D_h)U^{m+1}=(I+\tfrac\tau2\kappa D_h)U^m+\tfrac\tau2(f^{m+1}+f^m).
\]
对齐次边界的离散拉普拉斯本征模态，令 $z\ge0$ 为时间步长乘衰减率，两者放大因子是 $(1+z)^{-1}$ 和 $(1-z/2)/(1+z/2)$，模均不超过一。因此它们在离散 $L^2$ 意义下无上述时间步长稳定限制，但大步长仍降低精度；后一因子还可能为负，导致高频振荡。

\section{波动方程的离散传播}
标准中心格式为
\[
 U_j^{m+1}=2U_j^m-U_j^{m-1}+
 r^2(U_{j+1}^m-2U_j^m+U_{j-1}^m)+\tau^2 f_j^m,
 \qquad r=\frac{c\tau}{h}.
\]
设初速度采样为 $v_j$，可用 $U_j^1=g(x_j)+\tau v_j+\frac12\tau^2(c^2D_hg(x_j)+f_j^0)$ 启动。形式截断误差为 $O(\tau^2+h^2)$，前提是对应四阶导数有界。

代入频率模态，时间放大因子 $G$ 满足
\[
 G^2-2\bigl(1-2r^2\sin^2(\theta/2)\bigr)G+1=0.
\]
标准步长条件为 $r\le1$；若希望在一般稳定性论证中避开端点退化，可固定 $r\le r_0<1$。当 $r>1$，存在高频使一个根的模大于一。\emph{柯朗--弗里德里希斯--刘易条件（Courant--Friedrichs--Lewy condition，CFL condition）}要求数值依赖域能够覆盖物理依赖域，本例即 $h/\tau\ge c$。该几何条件一般是必要限制，不能脱离具体格式独自保证稳定。

在一步更新中，节点只使用左右各一个相邻点，所以经过 $m$ 步，最远只能使用距离 $mh$ 的初值；真实波在相同时间 $m\tau$ 内却能传播距离 $cm\tau$。若 $cm\tau>mh$，数值计算没有取得本应影响当前点的数据，不可能一致地再现任意初始扰动。这给出了 CFL 条件的几何来源。热方程的 $\tau=O(h^2)$ 限制则来自显式扩散权重及放大因子，不能简单套用波速解释。

\section{有限元与变分离散}
对于 $-u''=f$、零端点，在同一空间网格上取连续分片线性函数空间 $V_h\subset H_0^1(0,\ell)$，以节点帽函数 $\phi_i$ 为基。\emph{有限元法（finite element method）}求 $u_h=\sum_iU_i\phi_i$，使
\[
 \int u_h'v_h'=\int fv_h\quad(v_h\in V_h),\qquad
 KU=F,\quad K_{ij}=\int\phi_i'\phi_j',\quad F_i=\int f\phi_i.
\]
均匀网格上 $K_{ii}=2/h$、$K_{i,i+1}=K_{i+1,i}=-1/h$，其余非邻项为零。边界条件通过空间的选择实现，载荷积分的数值近似则另有求积误差。

\begin{example}[两项未知量的有限元方程]
对 $-u''=1$、$u(0)=u(1)=0$，把区间分成三段，节点为 $0,1/3,2/3,1$。每个内部帽函数的面积为 $1/3$，故
\[
 \begin{pmatrix}6&-3\\-3&6\end{pmatrix}
 \begin{pmatrix}U_1\\U_2\end{pmatrix}
 =\begin{pmatrix}1/3\\1/3\end{pmatrix},
 \qquad U_1=U_2=1/9.
\]
精确解 $u=x(1-x)/2$ 在两个内部节点的值也为 $1/9$，但单元内部的抛物线与分片直线仍不同，不能由节点相等便宣布整体误差为零。本例每个单元上的导数误差平方积分为 $h^3/12$，求和得到 $\|u'-u_h'\|_2=h/\sqrt{12}$，恰好展示了一阶的能量误差。
\end{example}

一般有界、强制双线性型下，精确弱解与伽辽金解满足\emph{伽辽金正交性（Galerkin orthogonality）}$a(u-u_h,v_h)=0$。记有界常数为 $M$、强制常数为 $\alpha$，则对任意 $v_h$，
\[
 \alpha\|u-u_h\|_V^2\le a(u-u_h,u-v_h)
 \le M\|u-u_h\|_V\|u-v_h\|_V.
\]
约去一项并取下确界，得到\emph{塞亚引理（C\'ea's lemma）}
\begin{equation}
 \|u-u_h\|_V\le\frac M\alpha\inf_{v_h\in V_h}\|u-v_h\|_V.
 \label{eq:cea}
\end{equation}
它把求解误差归结为逼近能力。例如一维 $u\in H^2$ 时，分片线性插值给出 $\|u-u_h\|_{H^1}\le Ch\|u\|_{H^2}$；更高维还需网格形状规则性，角点奇异会降低预期阶数。

\section{综合例题与学习检查}
\begin{example}[手算显式热格式的一步]
取 $\kappa=1$、$\ell=1$、$h=1/4$，零边界、零源项，并取内部温度向量 $(U_1^0,U_2^0,U_3^0)=(0,1,0)$。若 $\tau=1/64$，则 $\mu=1/4$，一步后得到 $(1/4,1/2,1/4)$，每个新值都是邻近旧值的非负加权平均。若改取 $\tau=1/16$，则 $\mu=1$，一步后得到 $(1,-1,1)$，已经从非负初值产生负温度；对应高频还会被反复放大。一个时间步的符号异常与频率稳定性分析在这里相互印证。
\end{example}

\begin{example}
为零边界热方程设计一个可以核对误差的计算实例。
\end{example}
\begin{solve}
在 $0<x<1$ 取已知函数 $u_*(x,t)=e^{-t}\sin(\pi x)$。由方程反算
\[
 f(x,t)=(\kappa\pi^2-1)e^{-t}\sin(\pi x),\qquad g(x)=\sin(\pi x).
\]
这组数据使 $u_*$ 成为精确解。使用式\eqref{eq:heat-explicit-scheme}，取固定 $0<\mu<1/2$ 即 $\tau=\mu h^2/\kappa$，在相同终止时间比较网格最大误差；当非零的二阶主误差项占主导时，网格步长减半后误差约缩小为四分之一。这是渐近预期，还需避免主误差项偶然抵消、时间终点不一致或舍入误差占主导的情形。
\end{solve}
\begin{enumerate}
 \item 实现或手工推进热方程前三个时间层，分别选择满足与违反稳定限制的参数。
 \item 对前述受迫弦振动例题比较模态精确解与中心差分结果，区分共振、相位误差及数值不稳定。
 \item 对 $-u''=1$ 的区间问题组装三对角有限元方程，用 $u=x(\ell-x)/2$ 检查误差。
\end{enumerate}
任何综合计算都应交代模型、数据、算法、网格与终止条件，并检查边界残差、误差或守恒量。图像的平滑性不足以判定数值正确性。

